\begin{figure*}[t]
\centering
\begin{tikzpicture}[x=1cm,y=1cm,
  stage/.style={draw=black!50,fill=black!6,rounded corners=1.5pt,align=center,font=\scriptsize,text width=2.5cm,minimum height=1.05cm,inner sep=2pt,anchor=north},
  ex/.style={draw=black!25,rounded corners=1.5pt,align=left,font=\scriptsize,text width=2.5cm,minimum height=2.45cm,inner sep=2.5pt,anchor=north},
  arr/.style={-{Latex[length=1.6mm]},black!60,line width=0.5pt}]
\node[stage] (s0) at (0.00,0) {\textbf{Agent READMEs dataset}\\\nPackageFiles{} agent files};
\node[stage] (s1) at (2.95,0) {\textbf{Same commit}\\README and CONTRIBUTING at the agent file's snapshot commit};
\node[stage] (s2) at (5.90,0) {\textbf{Filter and sample}\\\nEligible{} eligible files; \nPerAgent{} per file type; \nSampleRepos{} repositories};
\node[stage] (s3) at (8.85,0) {\textbf{Sections}\\split at headings; \nUnits{} sections};
\node[stage] (s4) at (11.80,0) {\textbf{Labelling (RQ1, RQ2)}\\up to 10 sections per repository; \nQueued{} labelled};
\node[stage] (s5) at (14.75,0) {\textbf{Detector (RQ3)}\\score of each section against the human sections};
\draw[arr] (s0.east) -- (s1.west);
\draw[arr] (s1.east) -- (s2.west);
\draw[arr] (s2.east) -- (s3.west);
\draw[arr] (s3.east) -- (s4.west);
\draw[arr] (s4.east) -- (s5.west);
\node[ex] (e0) at (0.00,-1.35) {\texttt{AGENTS.md} of the repository microsoft/\allowbreak generative-ai-with-javascript (995 stars).};
\draw[black!30,densely dotted,line width=0.5pt] (s0.south) -- (e0.north);
\node[ex] (e1) at (2.95,-1.35) {Snapshot commit \texttt{2e36ef3} (3 October 2025). Human documents: README.md and .github/\allowbreak CONTRIBUTING.md.};
\draw[black!30,densely dotted,line width=0.5pt] (s1.south) -- (e1.north);
\node[ex] (e2) at (5.90,-1.35) {Drawn for the \texttt{AGENTS.md} group; labelled in the main set, so it is part of the RQ3 test data.};
\draw[black!30,densely dotted,line width=0.5pt] (s2.south) -- (e2.north);
\node[ex] (e3) at (8.85,-1.35) {16 agent, 10 README and 6 CONTRIBUTING sections. This one: \emph{Local Development} (30 words, text below).};
\draw[black!30,densely dotted,line width=0.5pt] (s3.south) -- (e3.north);
\node[ex] (e4) at (11.80,-1.35) {\textbf{Coverage:} Agent-only. \textbf{Categories:} Configuration and Environment; Build and Run. \textbf{Need:} Development Environment and Tools. \textbf{Barrier:} Local environment setup hurdles.};
\draw[black!30,densely dotted,line width=0.5pt] (s4.south) -- (e4.north);
\node[ex] (e5) at (14.75,-1.35) {TF-IDF, 4 items; best matches 0.02, 0.03, 0.08, 0.03. Item mean $0.040 < 0.084$: agent-only (correct). Closest human section: \emph{Submitting a Pull Request (PR)}.};
\draw[black!30,densely dotted,line width=0.5pt] (s5.south) -- (e5.north);
\node[draw=black!25,fill=black!3,rounded corners=1.5pt,align=left,font=\scriptsize,text width=17.05cm,inner sep=3pt,anchor=north west] at (-1.25,-3.95) {\textbf{Example section} (\texttt{AGENTS.md}, \emph{Local Development}): 1. Ensure Node.js 20+ is installed 2. Set \texttt{GITHUB\_TOKEN} environment variable 3. Navigate to the directory you want to work with 4. Run \texttt{npm install} followed by \texttt{npm start}};
\end{tikzpicture}
\Description{Six method stages in a row, from the Agent READMEs dataset to the detector, with one real agent-file section traced below each stage.}
\caption{Overview of the pipeline. Top: the six stages. Bottom: one real section from our data traced through every stage, labelled agent-only by the annotator and predicted agent-only by the detector.}
\label{fig:overview}
\end{figure*}
